%% ============================================================
%% 分片 16：原文第 6 节「Loop and Line Operators」（原文 glshort.tex 起于第 4483 行）
%% 由 tools/harvmac2tex.py 从 arXiv 源码（hep-th/0604151v3）自动生成。
%% 只译英文散文；公式、编号、\label、\cite 一律不改。
%% ============================================================

\subsection{拓扑 't Hooft 算子}

\label{sec:tophooft}

我们已在 $t=\pm i$、即 $\Psi=\infty$ 处的拓扑场论中找到了 Wilson 算子。按 $S$-对偶 $S:\Psi\to -1/n_{\frak g}\Psi$，在 $\Psi=0$ 处必定存在对偶的算子。拓扑 Wilson 算子是由 $G$ 的一个表示 $R$ 的选择来分类的。因此，既然 $S$ 交换 $G$ 与 $^L\negthinspace G$，那么出现在 $\Psi=0$ 处 TQFT 中的对偶算子就必须由 $^L\negthinspace G$ 的一个表示 $^L\negthinspace R$ 来分类。

Wilson 算子就是统计力学中常说的「有序」（order）算子的一个例子。有序算子由某个经典表达式（例如联络绕一个圈的和乐）构造，随后被解释为量子算子，并作为一项因子纳入路径积分。

有序算子的对偶常常是「无序」（disorder）算子。无序算子是通过指定经典场所应具有的奇异性来定义的，并在该奇异性存在的情况下作路径积分（或量子化）。在经典层面，有序算子与无序算子看起来像是完全不同类型的对象。但在量子层面，在适当的情形下它们具有相同的形式性质，并且常常被对偶互换。在我们的问题中，Wilson 算子的对偶是 't Hooft 算子，它是无序算子的一个例子。\footnote{关于在 $G=U(1)$ 情形下 Wilson 算子的对偶确为 't Hooft 算子的一个显式论证，见 \cite{ref89}。}

要定义与一个圈 ${\cal S}\subset M$ 相伴的 Wilson 圈算子，我们需要 ${\cal S}$ 的一个定向。't Hooft 圈算子需要的则是 ${\cal S}$ 的法丛的一个定向。事实上，方便的做法是把 ${\cal S}$ 的一个小邻域等同于 ${\cal S}\times \Bbb{R}^3$，并记 $\epsilon_3$ 为 $\Bbb{R}^3$ 上的一个体积形式。't Hooft 算子的定义只依赖于法丛的定向，而不依赖于这些选择的细节。若 $r$ 表示到 ${\cal S}$ 的距离（在某个度规下），则我们也可以写 % 式 (6.6)
\setcounter{equation}{5}
\begin{equation}\label{eq:congo}
\epsilon_3=dr\,d{\rm
Vol},
\end{equation}
 其中 $d{\rm Vol}$ 是固定 $r$ 处的二维球面上的体积形式。

在本文中，当我们就 $U(N)$ 或 $SU(N)$ 规范理论中的 't Hooft 算子写出显式公式时，我们把规范场看作秩 $N$ 向量丛上的联络。这意味着 $A$ 与 $\phi$ 是取值为 $N\times N$ 反 Hermite 矩阵的 1-形式；当 $N=1$ 时，它们就简单地是虚 1-形式。

让我们先描述 $U(1)$ 规范理论中最基本的 't Hooft 算子。在 $\Bbb{R}^4=\Bbb{R}\times \Bbb{R}^3$ 上取坐标 $(s,\vec x)$。奇异的 Dirac 单极子是 Maxwell 方程的一个经典解，它在 $s$ 的平移下以及 $\vec x$ 的旋转下不变，并在由 $\vec x=0$ 定义的直线 $L$ 上有奇异性。若 $F=dA$ 是 $U(1)$ 联络 $A$ 的曲率，则 Maxwell 方程读作 $dF=d\star F=0$，而它可以在 $\Bbb{R}^3\backslash \{0\}$ 上（即在 $\Bbb{R}^3$ 中挖去点 $\vec x=0$ 之后）由下式求解 % 式 (6.7)
\setcounter{equation}{6}
\begin{equation}\label{eq:izo}
F={i\over 2}\star d\left({1\over |\vec
x|}\right).
\end{equation}
  使用 Hodge $\star$ 算子显式地显示出对 $L$ 的法丛定向的依赖。要得到 $\Bbb{R}^4\backslash L$ 上的解，只需利用投影 $\Bbb{R}^4\backslash L\to \Bbb{R}^3\backslash \{0\}$，把 $F$ 拉回到 $\Bbb{R}^4\backslash L$ 上，作为 Maxwell 方程的解。
 $F$ 已被归一化，使得若 $Y$ 是包围 $\vec x=0$ 处奇异性的二维球面，则 $\int_Y \,iF/2\pi =1$。因此存在一个酉线丛 ${\cal L}\to \Bbb{R}^4\backslash L$，它容许一个曲率为 $F$ 的联络。这就是奇异的 Dirac 单极子。它以 $\int_Yc_1({\cal L})=1$ 为特征。

阿贝尔规范理论中最基本的 't Hooft 算子的定义是：路径积分须在 $\Bbb{R}^4\backslash L$ 上的线丛的联络上进行，而这些联络在 $L$ 附近与奇异的 Dirac 单极子相符。这里没有必要只考虑直线。$L$ 可以换成一般四维流形 $M$ 中任意一条闭曲线 ${{\cal S}}$，或者（若 $M$ 是测地完备的）换成任意一条通向无穷远的曲线。在曲线 ${{\cal S}}$ 上插入 't Hooft 圈算子 $T_{(0)}$ 的路径积分，是通过对 $M\backslash {{\cal S}}$ 上的阿贝尔规范场\footnote{若 ${\cal S}$ 在 $H_1(M,\Bbb{Z})$ 中的同调类非零，则不存在这样的阿贝尔规范场，我们规定该路径积分为零。其对偶的说法是：若 ${\cal S}$ 的同调类非零，则 Wilson 算子 $W(R,{\cal S})$ 在对所有联络积分之后的期望值为零；这一为零来自用平坦线丛作扭化时的行为。}作路径积分而得到的，并要求曲率在 ${{\cal S}}$ 附近发散如 $F\sim {i\over 2}\star_3d (1/r)$，其中 $r=|\vec x|$ 现在是到 ${{\cal S}}$ 的距离，而 $\star_3$ 是作用在形式上的一个运算，它在 ${{\cal S}}$ 附近看起来就像法向于 ${{\cal S}}$ 的平面上的 $\star$ 算子。更一般地，取定一个整数 $m$ 并要求 $F$ 的奇异部分为 $F\sim{(im/ 2)}\star_3 d (1/r)$，就得到电荷为 $m$ 的 't Hooft 圈算子 $T_{(0)}(m)$。

现在让我们考虑规范群为 $U(1)$ 的 ${\EUN}=4$ 超对称杨–米尔斯理论。我们可以要求规范场具有 Dirac 单极子奇异性、而其余场非奇异，以此定义一个 't Hooft 算子。但这样的算子与 $\Psi=0$ 处的拓扑对称相容吗？

我们可以如下判定这个问题。要到达 $\Psi=0$，我们取 $t=1$ 且 $\tau$ 为虚数。为使 't Hooft 算子保持拓扑对称，场在 ${{\cal S}}$ 附近的奇异行为必须与该对称相容。由于问题只涉及 ${{\cal S}}$ 附近的局部行为，方便的做法是暂时回到 ${{\cal S}}$ 为 $\Bbb{R}^4$ 中直线 $L$ 的情形。

 一组场 $(A,\phi)$ 保持拓扑对称的条件已在式 \eqref{eq:loopme} 中求出。若这些方程与 't Hooft 算子所要求的那类奇异性相容，则 't Hooft 算子保持超对称。否则，在 't Hooft 算子存在时任何场组态都不可能超对称，我们就说 't Hooft 算子破坏超对称。

在 $t=1$ 且对阿贝尔规范理论，所需条件读作 % 式 (6.8)
\setcounter{equation}{7}
\begin{equation}\label{eq:nurot}
(F+d\phi)^+=(F-d\phi)^-=0.
\end{equation}
 若 $F$ 具有 Dirac 单极子奇异性而 $\phi$ 非奇异，这些方程就不被满足。要得到保持拓扑对称的 't Hooft 算子，正如 Wilson 的情形一样，我们必须以适当的方式把 $\phi$ 包含进来。

为使方程 \eqref{eq:nurot}\ 得到满足，$\phi$ 必须伴随 $F$ 的 Dirac 单极子奇异性而具有何种奇异性？这个问题有一个简单的答案：我们应取 $\phi=(i/2|\vec x|)ds$，其中 $s$ 是沿 $L$ 的「时间」坐标。于是，若取
% 式 (6.9)
\setcounter{equation}{8}
\begin{equation}\label{eq:hildo}
\begin{aligned}F& ={i\over 2}\star_3 d{1\over |\vec x|}\\
                    \phi& = {i\over 2|\vec x|} ds,\\\end{aligned}
\end{equation}

其中 $\star_3$ 是沿法向于 $L$ 的方向的三维 Hodge $\star$ 算子。

这就告诉我们，对于一般的闭一维流形 ${{\cal S}}\subset M$，如何定义在 $\Psi=0$ 处保持拓扑对称的基本 't Hooft 算子 $T$：要计算该算子的矩阵元，就对在 ${{\cal S}}$ 附近具有 \eqref{eq:hildo} 所描述奇异性的场作路径积分。类似地，对任意整数 $m$，要得到电荷为 $m$ 的拓扑 't Hooft 算子 $T(m)$，只需要求场沿 ${{\cal S}}$ 的奇异性是 \eqref{eq:hildo} 所描述奇异性的 $m$ 倍。

\par\medskip{\it 推广到任意 $G$ 与 Langlands 对偶}

下一步是把这一构造推广到任意 $G$。这可以简单地做到 \cite{ref90}。我们只需任取一个同态 $\rho:U(1)\to G$，然后要求场沿 ${{\cal S}}$ 的奇异性看起来像 \eqref{eq:hildo} 中奇异性在 $\rho$ 下的像，以此定义拓扑 't Hooft 算子。

这样定义的算子只依赖于 $\rho$ 的共轭类。但同态 $\rho:U(1)\to G$ 的共轭类由对偶群 $^L\negthinspace G$ 的最高权来分类。换言之，它们由 $^L\negthinspace G$ 的不可约表示来分类。

因此，对 $^L\negthinspace G$ 的每个不可约表示 $^L\negthinspace R$ 以及每个一维流形 ${{\cal S}}$，我们都得到了一个拓扑 't Hooft 算子 $T({}^L\negthinspace R,{{\cal S}})$。与 Wilson 的情形一样，若上下文清楚，我们就简单地称此算子为 $T({}^L\negthinspace R)$ 或 $T$。$G$ 的 't Hooft 算子由 $^L\negthinspace G$ 的表示来分类这一事实，就是经典的 GNO 对偶（GNO duality）\cite{ref10}，它近来被重新解释为关于算子而非态的对偶 \cite{ref90}。

't Hooft 算子的分类还与 Grothendieck 对 ${\Bbb {CP}}^1$ 上 $G$-丛的分类密切相关。按照这一分类，${\Bbb {CP}}^1$ 上的一个 $G$-丛通过一个同态 $\rho:U(1)\to G$（或其复化 $\rho:\Bbb{C}^*\to G_{\Bbb{C}}$）与基本线丛 ${\cal O}(1)\to\Bbb{CP}^1$ 相联系。我们把环绕 't Hooft 算子所定义奇异性的一个小二维球面与 $\Bbb{CP}^1$ 等同起来。黎曼面上的每个 $G$-丛，即便不稳定，也容许一个满足 Yang-Mills 方程的联络 \cite{ref91}；而在 $\Bbb{CP}^1$ 上，Yang-Mills 方程的每个解都等价于用来构造 't Hooft 算子的阿贝尔解之一。所以 't Hooft 算子的分类与 $\Bbb{CP}^1$ 上 $G$-丛的分类相同。

\par\medskip{\it 拓扑非平凡的 't Hooft 算子}

$S^2=\Bbb{CP}^1$ 上的全纯 $G$-丛在拓扑上可以非平凡，尽管 $G$-丛的拓扑分类远不如全纯分类精细。

$S^2$ 上的 $G$-丛由一个特征类 $\xi\in H^2(S^2,\pi_1(G))\cong \pi_1(G)$ 分类。但 $\pi_1(G)$ 自然同构于 $^L\negthinspace G$ 的中心。用于构造 't Hooft 算子 $T(^L\negthinspace R)$ 的 $S^2$ 上丛在拓扑上非平凡，当且仅当 $^L\negthinspace G$ 的中心在 $^L\negthinspace R$ 上非平凡地作用。

\par\medskip{\it 混合 Wilson-'t Hooft 算子：阿贝尔情形}

到目前为止，我们在 $\Psi=\infty$ 处找到了 Wilson 线算子，在 $\Psi=0$ 处找到了 't Hooft 线算子。通过施加对偶变换，我们找到 $\Psi$ 的其它特殊值，在这些值处存在保持拓扑对称的线算子。事实上，$\Psi$ 的所有有理值、并且只有这些值，可以由 $\Psi=0$（或 $\Psi=\infty$）通过对偶变换得到。无论对偶群 $\Gamma$ 是 $SL(2,\Bbb{Z})$（当 $G$ 为单边（simply-laced）型时就是如此）还是 $G$ 非单边时出现的某个 Hecke 群，这一点都成立。

要到达 $\Psi$ 的有理值，我们取 $t=1$ 且 $\theta=2\pi \Psi$。于是，我们在规范场的作用量中添加一项 % 式 (6.10)
\setcounter{equation}{9}
\begin{equation}\label{eq:into}
{{\cmmib I}}_\theta=i{\Psi\over 4\pi}\int
\Tr\,F\wedge F.
\end{equation}
 至此，有必要假定 $M$ 可定向，比如说通过选取一个体积形式 $\epsilon_4$。 ${\EUN}=4$ 超对称杨–米尔斯理论在不可定向的四维流形 $M$ 上也有意义。但我们此刻想考察其作用的相互作用 \eqref{eq:into}，只有在 $M$ 可定向时才能引入，因为定义积分 $\int_M\Tr\,F\wedge F$ 需要这一点。（更一般地，如我们在第~\ref{sec:canonpar}~节末尾所指出的，$\Psi$ 在定向反转下变号，所以要到达非零的有理 $\Psi$，我们需要 $M$ 上的一个定向。）因此我们假定 $M$ 可定向，并用体积形式 $\epsilon_4$ 的选择来表达这一方向。

定义支撑在 ${\cal S}$ 上的 Wilson 圈算子，需要用一个 1-形式 $ds$ 给出 ${\cal S}$ 的定向。定义支撑在 ${\cal S}$ 上的 't Hooft 圈算子，则需要 ${\cal S}$ 的法丛的一个定向 $\epsilon_3$。一旦 $M$ 可定向，我们就可以要求 ${\cal S}$ 的定向与其法丛的定向是相容的，即要求 % 式 (6.11)
\setcounter{equation}{10}
\begin{equation}\label{eq:delo}
\epsilon_4=ds\wedge \epsilon_3
\end{equation}
 沿 ${\cal S}$ 成立。
把 ${\cal S}$ 在 $M$ 中的一个小邻域等同于 ${\cal S}\times \Bbb{R}^3$，我们可以把 $ds$ 延拓到该邻域上，使得 \eqref{eq:delo}\ 在整个 ${\cal S}\times \Bbb{R}^3$ 上保持成立。然后如 \eqref{eq:congo} 那样写 $\epsilon_3=dr\wedge d{\rm Vol}$，我们就有 % 式 (6.12)
\setcounter{equation}{11}
\begin{equation}\label{eq:nelo}
\epsilon_4=ds\wedge dr\wedge d{\rm Vol}=-dr\wedge
ds\wedge d{\rm Vol}.
\end{equation}

在 $\Psi$ 的有理值处保持拓扑对称的线算子称为混合 Wilson-'t Hooft 算子。由于这些算子将兼具 't Hooft 算子与 Wilson 算子的性质，为定义它们，重要的是在 ${\cal S}$ 及其法丛上选取相容的定向，正如我们在上一段所做的那样。为构造这些算子，我们先考虑 $G=U(1)$ 的情形。

乍看起来，我们如此定义的 't Hooft 算子似乎对任意 $\Psi$ 都已经具有拓扑对称。但事实并非如此，因为 't Hooft 算子的奇异性使得 ${{\cmmib I}}_\theta$ 在拓扑对称下并不不变。

这需要一些解释。设 $A_0$ 是一个固定的联络，其曲率沿曲线 ${{\cal S}}$ 具有 Dirac 单极子奇异性。在电荷为 $m$ 的算子 $T_m({{\cal S}})$ 存在时的路径积分按如下方式计算：取规范场为 $A=mA_0+\hat A$，其中 $\hat A$ 在 ${{\cal S}}$ 附近光滑；然后在对 $\hat A$ 积分、$A_0$ 固定的情况下计算路径积分。$A_0$ 的精确选择无关紧要（只要它具有正确的奇异性），因为我们会对所有 $\hat A$ 积分。

作为朝这个方向的一步，让我们把 ${{\cmmib I}}_\theta$ 作为 $A$ 的函数来计算。我们写 $F=mF_0+\hat F$，其中 $F_0$ 与 $\hat F$ 分别是 $A_0$ 与 $\hat A$ 的曲率。由于我们将遇到一些奇异性，我们引入一个截断。设 $V_\epsilon$ 是以 ${{\cal S}}$ 为中心、半径为 $\epsilon$ 的管状邻域，我们将计算 $\int_{M-V_\epsilon} F\wedge F$，并考虑 $\epsilon\to 0$ 的极限。

其中一项来自 $\int_{M-V_\epsilon}m^2 F_0\wedge F_0$。这一项不太有意思，因为它与积分变量 $\hat A$ 无关。\footnote{这一项对 ${{\cmmib I}}_\theta$ 的贡献当 $\epsilon\to 0$ 时实际上收敛，因为 $F_0$ 的奇异部分秩为二。}另一项对我们当前的目的同样没有意思的贡献是 $\int_{M-V_\epsilon}\hat F\wedge \hat F$。由于 $\hat A$ 光滑，当 $\epsilon\to 0$ 时它有极限，即 $\int_M \hat F\wedge \hat F$。这个极限泛函是一个拓扑不变量，因此当然在拓扑对称所生成的 $\hat A$ 改变下不变。

有意思的贡献是交叉项 $\tilde {{\cmmib I}}_\theta =(im\Psi/2\pi)\int_{M-V_\epsilon}\,F_0\wedge \hat F$。我们想判定这个表达式是否被拓扑对称的生成元 $\delta_T$ 消灭。我们用 $\delta_T A_0=0$（因为 $A_0$ 是固定的规范场）与 $\delta_T\hat A=\delta_TA$ 来定义 $\delta_T$ 对 $A_0$ 与 $\hat A$ 的作用。我们有 $\delta_T \hat F=d (\delta_T\hat A)$，所以
% 式 (6.13)
\setcounter{equation}{12}
\begin{equation}\label{eq:juqu}
\delta_T\tilde {{\cmmib I}}_\theta={im\Psi\over
2\pi}\int_{M-V_{\epsilon}}d\left(F_0\wedge \delta_T\hat
A\right)={im\Psi\over 2\pi}\int_{\partial V_\epsilon} F_0\wedge
\delta_T\hat A.
\end{equation}
 要为这里的分部积分辩护，我们必须用 $M-V_\epsilon$ 的外法向量与 $M$ 的定向形式缩并来给出 $\partial V_\epsilon$ 的定向。由于 $M-V_\epsilon$ 的外法向是 $-d/dr$，我们从 \eqref{eq:nelo}\ 看出 $\partial V_\epsilon$ 必须由 $ds\wedge d{\rm Vol}$ 定向。

现在 $V_\epsilon$ 的边界是曲线 ${{\cal S}}$ 上的一个二维球丛。对 $V_\epsilon\to {{\cal S}}$ 的纤维上的积分可以利用 $F_0$ 在一条纤维上的积分为 $-2\pi i$ 这一事实来计算。于是我们得到
% 式 (6.14)
\setcounter{equation}{13}
\begin{equation}\label{eq:hugu}
\delta_T {{\cmmib I}}_\theta=m\Psi \int_{{\cal
S}}\delta_T \hat A.
\end{equation}

所以 ${{\cmmib I}}_\theta$ 在 't Hooft 算子存在时破坏拓扑对称；等价地说，当 $\theta\not= 0$ 时 't Hooft 算子破坏了拓扑对称。我们怎样才能恢复该对称？让我们回想量子力学中处理作用量 ${{\cmmib I}}$ 的办法：对所有场积分时带上因子 $\exp(-{{\cmmib I}})$。\footnote{在 Lorentz 号差下，这个因子应是 $\exp(i{{\cmmib I}})$，而 ${{{\cmmib I}}}_\theta$ 会多一个因子 $i$，得到的结果相同。}因此，把 \eqref{eq:hugu}\ 重新表述为 $\exp(-{{\cmmib I}}_\theta)$ 的变分公式是有用的：
% 式 (6.15)
\setcounter{equation}{14}
\begin{equation}\label{eq:jugu}
\delta_T(\exp(-{{\cmmib
I}}_\theta))=-m\Psi\,\exp(-{{\cmmib I}}_\theta)\, \int_{{\cal S}}
\delta_T\hat A.
\end{equation}

{}从这里出发，只差一小步就能看出
% 式 (6.16)
\setcounter{equation}{15}
\begin{equation}\label{eq:flugu}
\delta_T\left(\exp\left(m\Psi\int_{{\cal S}} \hat
A\right)\exp\left(-{{\cmmib I}}_\theta\right)\right)=0.
\end{equation}
 因此，只要在路径积分中把 Wilson 算子 $\exp(m\Psi\int_{{\cal S}} \hat A)$ 作为额外的因子乘进去，我们就能恢复拓扑对称。 （注意这里必须使用朴素的 Wilson 算子，而不是带 $\phi$ 贡献的超对称 Wilson 算子。）

然而，对一般的实数 $n$，表达式 $\exp(-n\int_{{\cal S}} \hat A)$ 并非规范不变。它规范不变当且仅当 $n$ 是整数。所以要恢复拓扑对称，$n=-m\Psi$ 必须是整数。换句话说，$\Psi$ 必须是有理数 $-n/m$。对有理的 $\Psi=-n/m$，我们已看到拓扑对称被电荷为 $n$ 的 Wilson 算子 $\exp(-n\int_S \hat A)$ 与电荷为 $m$ 的 't Hooft 算子的乘积所保持。该乘积称为 Wilson-'t Hooft 算子。我们把它记为 $WT_{n,m}({{\cal S}})$。

对有理的 $\Psi$，我们把 $\Psi$ 写成 $\Psi=-n/m$，其中 $n$ 与 $m$ 是互素的整数。在对偶变换 $\Psi\to (a\Psi+b)/(c\Psi+d)$ 下，$n$ 与 $m$ 按
% [WARN] 公式 \eqn\eqref{eq:yto} 无编号或重复
\begin{equation*}
\tag*{(6.17)}
\left(\begin{matrix} m  & n \\\end{matrix}\right)\to  \left(\begin{matrix} m &
n\\\end{matrix}\right) \left(\begin{matrix} d & -b\\ -c & a \\\end{matrix}\right) .
\end{equation*}
 变换。正如式 \eqref{eq:xilj} 中那样，这是磁荷与电荷 $\eurm m$ 与 $\eurm n$ 在对偶下的标准变换。其原因是：在静态 Wilson-'t Hooft 算子 $WT_{n,m}$（且没有其它电荷或源）存在时，把阿贝尔规范理论量子化所得的态具有 $({\eurm m}~{\eurm n})=(m~n)$。对磁荷，这一点相当清楚，因为 $m$ 与 ${\eurm m}$ 都是通过第一陈类定义的（分别在无穷远处或在奇异性附近测量）。对电荷，可以通过求解 Maxwell 方程、确定 Wilson 算子存在时电场的渐近行为来证明。我们刚刚完成的计算表明，$n=-\Psi m$ 在 $\Psi\to\Psi+1$ 下的变换与 \eqref{eq:yto} 中所断言的一致。

\par\medskip{\it 混合 Wilson-'t Hooft 算子：非阿贝尔情形}

与 't Hooft 算子一样，Wilson-'t Hooft 算子到一般单规范群 $G$ 的推广，是通过同态 $\rho:U(1)\to G$ 把阿贝尔构造嵌入 $G$ 中来完成的。然而，这需要一些解释。

我们写 $A=\rho(A_0)+\hat A$，其中 $A_0$ 是沿 ${{\cal S}}$ 具有 Dirac 奇异性的奇异 $U(1)$ 规范场，而 $\hat A$ 光滑。（对 $G=U(1)$，同态 $\rho$ 在李代数层面就是乘以某个整数 $m$，所以我们写 $A=mA_0+\hat A$。）然而现在我们必须澄清 $\hat A$ 光滑意味着什么。由于曲率 $F=dA+A\wedge A$ 中有非线性项，我们必须要求
% 式 (6.18)
\setcounter{equation}{17}
\begin{equation}\label{eq:milfo}
[\rho(A_0),\hat A]=0
\end{equation}
 沿 ${{\cal S}}$ 成立，否则曲率沿 ${{\cal S}}$ 的奇异部分将依赖于 $\hat A$，不与 Dirac 奇异性相符。

为使条件 \eqref{eq:milfo}\ 有意义，我们还必须把沿 ${{\cal S}}$ 的规范变换限制到与 $\rho(U(1))$ 交换的子群 $H\subset G$。于是，限制到 ${{\cal S}}$ 上时，$\hat A$ 的结构群是 $H$。

$H$ 的李代数有一个 $H$-不变投影 $\pi_\rho$ 到 $\rho(\frak{u}(1))$，即 $\rho(U(1))$ 的李代数。沿 ${{\cal S}}$，我们可以定义投影 $\pi_\rho(\hat A)$，它是一个取值于 $\rho(\frak{u}(1))$ 的规范场。

与阿贝尔情形相同的计算表明，对应于 $\rho$ 的 't Hooft 算子当 $\Psi\not= 0$ 时不再在拓扑对称下不变。取而代之的是
% 式 (6.19)
\setcounter{equation}{18}
\begin{equation}\label{eq:yugo}
\delta_T {{\cmmib I}}_\theta={\Psi}\int_{{\cal S}}
\Tr\,\rho(1)\,\delta_T\hat A.
\end{equation}
 这里 $\rho(1)$ 就是 $1\in \frak u(1)$ 在李代数同态 $\rho:\frak u(1)\to \frak g$ 下的像。表达式 $\Tr\,\rho(1)\,\delta_T \hat A$ 是 $\pi_\rho(\delta_T\hat A)$ 的一个倍数。

与阿贝尔情形一样，只要把 't Hooft 算子乘上
 % 式 (6.20)
\setcounter{equation}{19}
\begin{equation}\label{eq:yrot}
\exp\left(\Psi\int_{{\cal S}} \Tr \,\rho(1)\hat A\right).
\end{equation}
 就能恢复拓扑不变性。其规范不变性对 $\Psi$ 提出一个条件，非正式地说，就是 $\rho$ 必须能被 $\Psi$ 的分母整除。对任意有理的 $\Psi$，只要取合适的 $\rho$，这一条件就得到满足。

这样得到的算子称为 Wilson-'t Hooft 算子。底层的超对称规范理论有更一般的 Wilson-'t Hooft 算子，其电荷权与磁荷权线性无关 \cite{ref90}。拓扑场论中的 Wilson-'t Hooft 算子全部由 $\Psi=\infty$ 处的 Wilson 算子经对偶产生，所以它们的电荷权与磁荷权成正比。

关于 Wilson-'t Hooft 算子的近期相关讨论，另见 \cite{ref92}。

\begin{yizhu}[易错点]
Wilson 圈算子需要 ${\cal S}$ 的切向定向（1-形式 $ds$），而 't Hooft 圈算子需要的是法丛的定向（$\epsilon_3$）。两者并不独立：一旦 $M$ 可定向，\eqref{eq:delo} 的相容条件 $\epsilon_4=ds\wedge\epsilon_3$ 把二者联系起来。这一步在 \eqref{eq:juqu} 的分部积分中至关重要——正是用它才断定 $\partial V_\epsilon$ 的定向为 $ds\wedge d{\rm Vol}$，从而边界项的符号正确。若混用定向，会得到 $\delta_T{{\cmmib I}}_\theta=-m\Psi\int\delta_T\hat A$ 之类的符号错误。
\end{yizhu}

\begin{yizhu}[易错点]
式 (6.17) 中的 $(m\ n)$ 是行向量，其分量依次为磁荷与电荷，右乘的 $2\times 2$ 矩阵属于对偶群。原文的 harvmac 源把 2$\times$2 矩阵压成了 1$\times$2 的写法（矩阵中多出的空行即由此而来），阅读时须按「行向量右乘矩阵」理解，而不是把它当作标量等式。
\end{yizhu}

\begin{yizhu}[物理图像]
't Hooft 算子的分类等于 $\Bbb{CP}^1$ 上 $G$-丛的全纯分类：绕奇异点的二维球面把 $U(1)$ 单极子按同态 $\rho:U(1)\to G$ 嵌入 $G$，而 $\rho$ 的共轭类又对偶于 $^L\negthinspace G$ 的最高权，即 $^L\negthinspace G$ 的不可约表示。丛的拓扑类型由 $\pi_1(G)\cong Z(^L\negthinspace G)$ 刻画，故「拓扑非平凡」当且仅当 $^L\negthinspace G$ 的中心在相应的表示 $^L\negthinspace R$ 上非平凡地作用。这三层（同态、最高权、中心作用）是同一件事的三种说法。
\end{yizhu}

\begin{yizhu}[含糊点拨]
为什么只有有理的 $\Psi$ 才有保持拓扑对称的线算子？'t Hooft 算子使 $\theta$ 项在奇异背景下不再规范不变，其变化正比于 $m\Psi\int_{{\cal S}}\delta_T\hat A$；要抵消它必须乘上一个 Wilson 因子 $\exp(-n\int_{{\cal S}}\hat A)$，而后者规范不变当且仅当 $n$ 为整数。于是 $\Psi=-n/m$ 必须是有理数，且由 $\Psi=0$ 处的 't Hooft 算子与 $\Psi=\infty$ 处的 Wilson 算子按 $m:n$ 的比例混合而成。反过来，任何有理 $\Psi$ 都可由对偶变换从 $0$ 或 $\infty$ 得到。
\end{yizhu}

\begin{yizhu}[易错点]
非阿贝尔情形不能简单地把 $A=mA_0+\hat A$ 搬过来：曲率的非线性项 $A\wedge A$ 含有 $\rho(A_0)\wedge\hat A$，所以必须附加条件 \eqref{eq:milfo} $[\rho(A_0),\hat A]=0$，否则奇异部分会随 $\hat A$ 改变。相应地，沿 ${{\cal S}}$ 的规范变换只能取与 $\rho(U(1))$ 交换的子群 $H$，而 $\hat A$ 的结构群降为 $H$；把 $\hat A$ 投到 $\rho(\frak{u}(1))$ 上的 $H$-不变投影 $\pi_\rho$ 正是用来把非阿贝尔情形化归为阿贝尔情形的工具。
\end{yizhu}

\begin{yizhu}[物理图像]
式 \eqref{eq:juqu}--\eqref{eq:flugu} 是典型的「反常流入」式计算：$\theta$ 项在 't Hooft 奇异背景下的变分不是零，而是化为管状邻域边界上的通量积分；把 $F_0$ 在纤维上的积分 $\int_{\rm fiber}F_0=-2\pi i$ 代入，即得 $\delta_T{{\cmmib I}}_\theta=m\Psi\int_{{\cal S}}\delta_T\hat A$。再把这个变分写成 $\exp(-{{\cmmib I}}_\theta)$ 的变分，就能看出乘上适当的 Wilson 因子恰好抵消它。注意在 Lorentz 号差下插入因子是 $\exp(i{{\cmmib I}})$，结论不变。
\end{yizhu}

\begin{yizhu}[原文疑误]
原文第 59 行作 「This means that $A$ and $\phi$ are are one-forms valued in ...」，其中「are are」是重复词，应为单个 are（底本笔误）。译文按单次出现处理，未改动任何数学内容；此处不涉及公式。另原文在本小节两处重复声明「assume $M$ to be oriented」，属行文重复而非错误，译文照原文的重复结构分别译出。
\end{yizhu}
